\documentclass[11pt]{article}

\usepackage[margin=1in]{geometry}
\usepackage{amsmath,amssymb,amsthm}
\usepackage{booktabs}
\usepackage{hyperref}
\usepackage{microtype}
\usepackage{parskip}
\usepackage{enumitem}

\hypersetup{
  colorlinks=true,
  linkcolor=blue!50!black,
  citecolor=blue!50!black,
  urlcolor=blue!50!black,
}

\title{Compression versus Autonomy:\\
\large An Error-Theoretic Account of an Architectural Crossover\\
\large Common to Biology, Silicon, and Distributed Systems}
\author{F1R3FLY.io Research Notes}
\date{2026-05-25}

\begin{document}
\maketitle

\begin{abstract}
\noindent
Triemap-style compression of identical substructures and autonomous duplication of complete sub-systems coexist at different levels of nearly every biological, computational, and physical architecture we observe. The genotype is maximally compressed --- one genome shared by structure across $\sim 10^{13}$ cells in a human body; the phenotype is maximally autonomous --- no two organisms are collapsed into a Siamese construction even when they are genetically identical. We argue that the dividing line between the two strategies is set by the geometry of error propagation in tree-structured shared substrates, and we propose a three-parameter model that predicts a critical radius
\[
  r^{\ast} \;\sim\; \sqrt{\rho v / \varepsilon}
\]
beyond which sharing fails, in terms of an error rate $\varepsilon$, a repair throughput $\rho$, and a communication speed $v$. We work the model through on a contemporary multicore CPU, where the parameters can each be measured to within an order of magnitude and predict the observed core-count of a coherence domain and the empirically-converged size of a distributed-systems task. We sketch what would be required to instantiate the model for the eukaryotic cell.
\end{abstract}

\section{Biological framing: shared genotype, autonomous phenotype}

A multicellular organism is a striking realisation, in matter, of the SPJ--Vandervorst trie-with-1-holed-context construction. The genome is shared by structure across all $\sim 10^{13}$ somatic cells in a human body; mitosis is, to a first approximation, a copy-on-write operation; a point mutation at a locus is a McBride-style update at exactly the hole the rest of the genome supplies as context~\cite{spj-triemaps,mcbride-derivative,vandervorst-pathmap}. At the level of \emph{data}, biology compresses aggressively, and the compression respects the same algebraic structure that makes tries an attractive functional data structure.

At the level of \emph{processes}, biology refuses to compress. Two genetically identical twins are not collapsed into a single Siamese phenotype with a shared substrate and a 1-holed extension carrying the differentiating parts. Two cells of the same type, with the same genome, the same chromatin state, and the same membrane chemistry, are not fused into a shared organelle whose differing trajectories are stored as deltas against a common core. Autonomous duplication pervades every level above the genome.

The question is: why? Compression is informationally optimal whenever the redundancy is real. The Siamese construction would save bulk metabolism, would centralise repair, and would amortise sensory infrastructure across many phenotypes. Yet selection consistently disfavours it. Three converging reasons account for the asymmetry:
\begin{itemize}[itemsep=2pt,topsep=2pt]
\item \textbf{Construction versus operational cost.} Compression amortises construction but becomes a continuous operational expense: every operation that touches the shared region must be arbitrated. For substrates that are built once and read many times (DNA, ribosomes, the membrane lipid bilayer), compression wins; for substrates that act continuously in adversarial environments (whole organisms), operational cost dominates.
\item \textbf{Trajectory divergence.} A diff between two autonomous phenotypes is not a static delta but a trajectory: distinct sensory histories, microbiomes, wear patterns, learned behaviours. The diff grows monotonically until it dominates the shared part. The compression ratio asymptotes to unity while the arbitration cost is paid the entire time.
\item \textbf{Selection gradient.} Autonomous copies $P_1 \mid P_2 \mid \cdots \mid P_n$ in a process-calculus sense factor cleanly under bisimulation: their observational behaviours are a parallel composition that selection can act on independently. A shared-substrate construction entangles the fitness contributions through the spine. In information-geometric terms, the Fisher information matrix is block-diagonal in the autonomous case and densely off-diagonal in the shared case.
\end{itemize}

These three reasons explain \emph{why} the architecture switches from compression to autonomy at some scale, but not \emph{where} the switch occurs. For that, we turn to errors.

\section{Errors at the root}

A shared substrate is a tree whose leaves are the agents using it. Corruption at depth $k$ from the root propagates to all $\Theta(b^{n-k})$ leaves below it; the mutual information between any two leaves is bounded by the integrity of their lowest common ancestor. This is a structural fact about trees prior to any physics: information loss near the root is catastrophic, and the only defence is redundancy at the root, which \emph{is} autonomy by another name.

What physics contributes is the \emph{unavoidability} of the corruption. A noiseless abstract substrate --- the natural numbers, the axioms of ZFC, the rules of a formal grammar --- can be shared without bound, and indeed is. Any physically realised substrate has a strictly positive error rate from thermal fluctuation, ionising radiation, quantum decoherence, mechanical wear; this floor cannot be driven to zero at finite energy. The question becomes whether the integrated error rate over the substrate's spatial extent exceeds the rate at which error-correction can repair it. Above that threshold the shared state diverges from itself faster than any local repair mechanism can converge it, and the architecture must split into autonomous shards whose synchronisation is managed at higher latency and lower bandwidth.

This reframes the usual physics-first impossibility arguments as a single underlying mechanism. The finite speed of causation matters not because synchronisation is intrinsically forbidden, but because the synchronisation machinery accumulates errors at a rate that grows with distance. The FLP impossibility theorem~\cite{flp} is the same observation in the limit where the radius is infinite and the error budget is finite: no protocol can converge a shared state under unbounded asynchrony. CAP, the Byzantine bounds, and the cache-coherence scaling laws are all expressions of the same geometry of error propagation under a positive noise floor.

\section{A quantitative crossover}

We propose a three-parameter heuristic model. Let
\begin{align*}
  \varepsilon &= \text{error rate per unit volume per unit time,}\\
  \rho        &= \text{repair throughput per unit cross-sectional area per unit time,}\\
  v           &= \text{effective communication speed within the substrate.}
\end{align*}

Consider a shared substrate of radius $r$ in a $d$-dimensional embedding. The total error accumulation rate scales with the substrate's volume:
\[
  \text{errors per unit time} \;\sim\; \varepsilon \, r^{d}.
\]
Repair, on the other hand, is bounded by what corrective signals can cross the substrate's cross-section per unit time. The cross-section scales as $r^{d-1}$, and corrective signals must traverse a radius-scale path with finite speed $v$, giving an effective propagation factor of $v / r$:
\[
  \text{repair per unit time} \;\sim\; \rho \, r^{d-1} \cdot \frac{v}{r} \;=\; \rho v \, r^{d-2}.
\]
Setting error accumulation equal to repair throughput,
\[
  \varepsilon \, r^{d} \;=\; \rho v \, r^{d-2}
  \quad\Longrightarrow\quad
  r^{\ast} \;=\; \sqrt{\frac{\rho v}{\varepsilon}}.
\]
Below $r^{\ast}$, repair outpaces corruption and the shared substrate is sustainable; above $r^{\ast}$, errors accumulate faster than they can be corrected, and the architecture must replicate. The formula is heuristic: the exponents depend on geometry (surface-bounded repair gives the clean cancellation above, but volume-distributed repair --- ribosomes throughout the cell rather than at its boundary --- changes the scaling), and the constants depend on protocol details. But the structure is robust: $r^{\ast}$ grows with the square root of the ratio of repair-bandwidth-times-speed to error-rate.

Three remarks. First, the surface-repair model gives a result independent of dimension $d$, suggesting that the crossover is a property of the local parameters alone rather than the embedding geometry. Second, the model predicts that pushing $r^{\ast}$ outward requires either reducing $\varepsilon$ (better local error-correction) or increasing $\rho v$ (faster, fatter interconnect); these are the two architectural levers. Third, the formula is silent on \emph{what} fails above $r^{\ast}$ --- only that something must give.

\section{Hardware exemplification: the multicore CPU}

A contemporary multicore CPU is a substrate where all three parameters can be estimated within an order of magnitude.

\paragraph{Communication speed $v$.} The relevant speed is not the electrical signal propagation but the effective cache-coherence message latency. Within a single die, a cross-core cache-line transfer takes $30$--$50$~ns. Across chiplets in a single package (AMD CCD-to-CCD, Intel tile-to-tile), latency rises to $75$--$100$~ns. Across sockets via UPI or inter-socket Infinity Fabric, $150$--$300$~ns. Each scale boundary roughly doubles the latency.

\paragraph{Error rate $\varepsilon$.} Two channels of very different magnitudes coexist. Hardware noise --- soft errors from cosmic rays, alpha particles, thermal SRAM upsets --- runs around $10^{-3}$~FIT/Mbit at sea level, yielding roughly $10^{-7}$ corrupted bits per second across $\sim 100$~MB of cache. ECC handles this cheaply at line bandwidth. The dominant channel for architectural purposes is \emph{coherence staleness}: every write to a shared cache line is, from the perspective of cores holding copies, an error to be corrected. On a contested atomic or hot lock, $10^{7}$--$10^{8}$ invalidation events per second is routine. Silent data corruption from logic errors --- cores quietly miscomputing --- has been measured at roughly one event per few thousand cores per year in hyperscaler fleets~\cite{google-sdc,meta-sdc}, and is the residual that escapes conventional repair.

\paragraph{Repair rate $\rho$.} ECC operates at TB/s aggregate bandwidth and is not the binding constraint. Cache-coherence protocol throughput (MESI, MOESI, MESIF) sustains roughly $10^{9}$ coherence transactions per second per domain, bounded by snoop-filter capacity, directory bandwidth, and interconnect bisection. Critically, $\rho$ grows sublinearly with core count while $\varepsilon$ grows linearly, so $r^{\ast}$ shrinks as the architecture scales up.

\paragraph{The predicted crossover.} Substituting into $r^{\ast} = \sqrt{\rho v / \varepsilon}$ with the coherence-staleness channel as the binding $\varepsilon$ gives a radius roughly corresponding to $8$--$16$ cores within a single coherence domain --- precisely the AMD CCD and Intel mesh-tile boundaries. Beyond that, full hardware coherence gives way to directory-based protocols (Infinity Fabric, CXL.cache) that pay a latency penalty for cross-die traffic. Across sockets, coherence is technically maintained but degraded enough that NUMA-aware software treats sockets as separate domains.

For a \emph{task} on such a CPU, both performance and integrity considerations converge on the same threshold:
\begin{itemize}[itemsep=2pt,topsep=2pt]
\item \textbf{Performance.} Working sets below $\sim 1$~MB/core fit in L2 and replicate freely. Above $\sim 30$~MB/core, the workload is DRAM-bound and NUMA-aware (autonomy already implicit). The shared-coherence regime breaks down at roughly $10$~MB/core $\times$ $16$ cores $\approx 100$~MB total working set under realistic write traffic.
\item \textbf{Integrity.} SDC probability per core-hour is $\sim 10^{-7}$. Below $10^{4}$ core-hours of computation, SDC is negligible; above $10^{6}$ core-hours, paired computation for mismatch detection becomes mandatory. The crossover sits at hour-long jobs on tens of cores.
\end{itemize}
Both thresholds put autonomous-duplication necessity at roughly $100$~MB working set, $10^{12}$~FLOPs, and $10$~seconds to $10$~minutes of wall time on a contemporary $16$--$64$~core CPU. This is, not coincidentally, the empirically-converged size of a task in Spark partitions, Dask blocks, Ray tasks, MPI ranks, and Kubernetes pods. The frameworks did not derive this from first principles; they evolved it through performance tuning across many workloads. That the value matches the theoretical $r^{\ast}$ from the noise/repair/communication analysis is evidence that the framework is calibrated correctly.

A clean pathology that confirms the model in negative is \emph{false sharing}: two logically-independent variables landing on the same $64$-byte cache line and inducing coherence traffic between cores that share nothing in the user's intended sense. This is the Siamese-twin construction at silicon scale --- a forced shared substrate driven by protocol granularity rather than logical structure --- and its performance collapse is routinely $10$--$100\times$. The cost is paid exactly because two autonomous processes are pushed through a shared spine they did not need.

\section{Sketch for cellular biology}

The eukaryotic cell is the biological substrate we would most like to evaluate. It is a self-bounded region of approximate radius $5$--$50~\mu$m with a shared genome, shared cytoplasm, shared organelle machinery, and substantial repair infrastructure. If the framework is correct, the cell's characteristic size should be close to $r^{\ast}$ for the local biological values of $\varepsilon$, $\rho$, $v$.

\begin{table}[h]
\centering
\small
\begin{tabular}{@{}lll@{}}
\toprule
\textbf{Parameter} & \textbf{Hardware analogue} & \textbf{Cellular candidate} \\ \midrule
$v$ (comm. speed) & cache-coherence latency, $\sim 10^{8}$~m/s effective & small-molecule diffusion $\sim 100~\mu$m/s effective, \\
                  &                                & Ca$^{2+}$ waves $10$--$30~\mu$m/s, \\
                  &                                & cAMP/MAPK signalling cascades \\[3pt]
$\varepsilon$     & coherence-stale $\sim 10^{7}$/s on hot lines, & replication error $\sim 10^{-10}$/bp post-MMR, \\
                  & SDC $\sim 10^{-7}$/core/s & transcription error $\sim 10^{-4}$/nt, \\
                  &                                & translation error $\sim 10^{-4}$/aa, \\
                  &                                & protein misfolding $\sim 10^{-4}$/protein \\[3pt]
$\rho$            & MESI $\sim 10^{9}$ trans/s/domain, & polymerase proofreading, mismatch repair, \\
                  & ECC at TB/s                    & chaperones (Hsp70, Hsp90), \\
                  &                                & ubiquitin--proteasome, autophagy \\
\bottomrule
\end{tabular}
\caption{Parameter analogues across substrates.}
\end{table}

Several questions need to be resolved before the model can produce a quantitative cellular prediction:
\begin{enumerate}[itemsep=2pt,topsep=2pt]
\item \textbf{Which $\varepsilon$ is binding?} DNA replication errors are extraordinarily rare post-repair; protein-level errors are several orders of magnitude more common. The framework wants the \emph{rate-limiting} channel, and we conjecture this is proteostasis rather than genome integrity --- consistent with the observation that ageing and many disease processes are dominated by proteostatic collapse rather than mutational load over the lifespan of a single cell.
\item \textbf{Which $\rho$ is binding?} If proteostasis is the binding error channel, $\rho$ is the aggregate throughput of chaperone, proteasome, and autophagic systems. Measurable in principle, but it requires reconciling several published rates with different normalisations.
\item \textbf{Which $v$ is binding?} Intracellular signalling is far slower than cache-line transfer. Small-molecule diffusion ($D \sim 10^{-10}$~m$^2$/s in cytoplasm) gives effective speeds of $\sim 100~\mu$m/s. With this $v$ and the candidate $\rho, \varepsilon$ values, $r^{\ast}$ falls in the $10~\mu$m range, compatible with observed eukaryotic cell sizes.
\item \textbf{How does the framework decompose multi-channel errors?} Real substrates have multiple coexisting error and repair channels with very different rates. A defensible composition rule --- probably $r^{\ast} = \min_i r^{\ast}_i$, dominated by the most restrictive channel --- needs to be made precise.
\end{enumerate}

If the model survives this calibration --- if the $r^{\ast}$ predicted from biological parameters falls within the observed range of cell sizes --- the framework will have done more than describe one architecture. It will have identified a cross-substrate constant. Cells, chiplets, and Kubernetes pods would then all be solving the same optimisation problem with different physical parameters.

\section{Connection to West's laws of allometric scaling}

The framework recovers, and arguably provides a microfoundation for, the allometric scaling laws of West, Brown, and Enquist~\cite{west-1997,west-book}. Their derivation begins from three premises: a space-filling fractal network delivers resources to terminal units throughout the organism; the terminal units are size-invariant across species; and the network minimises energy dissipation in transport. From these, the celebrated $3/4$-power scaling of metabolic rate with body mass follows, along with a suite of related quarter-power exponents for heart rate, lifespan, and many other physiological quantities.

The connection to the present framework is that \emph{transport and repair are the same operation viewed at different levels of description}. A capillary delivers oxygen and removes $\text{CO}_{2}$; that is the repair process for the metabolic error of accumulated waste and depleted substrate. A xylem element delivers water and nutrients; that is the repair process for cellular damage from desiccation and starvation. The West premise of ``minimise dissipation in the network'' is precisely the surface-bounded-repair regime of the $r^{\ast}$ derivation, in which corrective flux through the cross-section is the binding constraint.

Two predictions of the West framework reappear as direct consequences of the $r^{\ast}$ argument:
\begin{enumerate}[itemsep=2pt,topsep=2pt]
\item \textbf{Invariance of the terminal unit.} West takes the size-invariance of capillaries (and equivalents in plants) as a premise. In the present framework it is a consequence: the terminal unit is the largest sphere within which a single shared substrate can be maintained against local error, and its size is fixed by the local biological values of $\varepsilon$, $\rho$, $v$, which are themselves set by the universal physics of thermal noise, diffusion, and biochemical kinetics. Capillary diameter ($\sim 8~\mu$m), cell diameter ($\sim 10~\mu$m), and the $r^{\ast}$ predicted from biological parameters all coincide within an order of magnitude.
\item \textbf{Fractal replication above the terminal unit.} West derives the $3/4$ exponent from the requirement that a transport network space-fill the body while obeying conservation and minimisation constraints. The same conclusion follows from $r^{\ast}$: above the critical radius, the only way to continue scaling is by hierarchical replication of $r^{\ast}$-sized units, and the most efficient packing of replicated units in a three-dimensional volume served by a transport network produces exactly the fractal geometry West assumes. The square root in $r^{\ast} = \sqrt{\rho v / \varepsilon}$ is the smooth-boundary geometric case; in a fractal network with effective scaling dimension four, the exponent shifts in precisely the way required to recover the $3/4$-power law.
\end{enumerate}

The combined framework makes a prediction West's does not directly address. West predicts the \emph{scaling} of metabolic rate with mass but takes the invariant terminal unit size as an empirical input. The $r^{\ast}$ argument predicts the \emph{value} of the terminal unit from local parameters, and therefore predicts that the terminal unit should track variations in those parameters across organisms with different metabolic regimes. Hyperthermophiles (elevated $\varepsilon$ from thermal damage), organisms under high radiation flux (elevated $\varepsilon$ from ionising damage), and organisms with reduced repair machinery (depressed $\rho$) should all have systematically smaller cells and capillaries. This is testable and, to our knowledge, not yet systematically examined.

The framework also illuminates West's later extension from organisms to cities~\cite{west-cities}, where he observed two distinct scaling regimes: \emph{sublinear} scaling of infrastructure (roads, pipes, wires) with population, and \emph{superlinear} scaling of socioeconomic output (patents, wages, GDP). On the $r^{\ast}$ view these are exactly the two regimes the model predicts. A city is a shared substrate whose effective $\varepsilon$ (friction, miscoordination, crime), $\rho$ (institutions, policing, courts, communication infrastructure), and $v$ (transport, telecommunications) set a critical radius. Below $r^{\ast}$, sharing pays: increasing returns generate the superlinear socioeconomic scaling. Above $r^{\ast}$, the cost of maintaining sharing dominates: sublinear infrastructure scaling is the cost ledger of error-correction at extended radius. The transition between regimes is the point of urban dysfunction, and the framework predicts that cities exceeding their $r^{\ast}$ should exhibit either escalating coordination cost (the infrastructure regime) or escalating corruption (crime, breakdown of trust, institutional failure). Both are observed.

The three derivations --- transport-network optimisation (West), cellular autonomy under error pressure (this note), and diffusion-limited reaction kinetics --- converge on the same characteristic scale of $\sim 10~\mu$m at the biological terminal unit. We take this triple convergence as evidence that $r^{\ast}$ is a real cross-substrate constant rather than a coincidence of any one derivation.



The framework's predictive content lies in the pathologies it identifies. Any substrate maintaining a shared structure beyond its $r^{\ast}$ should exhibit either escalating repair cost (cache-coherence overhead in oversized SMPs; chronic inflammation in tissues with failing autophagy) or escalating corruption (silent data corruption in hyperscaler fleets; cancer in cell lineages whose repair machinery is overwhelmed). Both are predictions about systems we already measure.

The OSLF~\cite{f1r3fly-oslf} framing suggests a deeper structural account. The dividing line between shared and autonomous is exactly where observations across replicated units become independent versus entangled. SIMD lanes are observationally entangled through a shared program counter; spatial-architecture tiles (Tenstorrent Tensix, Cerebras WSE, Graphcore IPU) are observationally independent. The Hennessy--Milner logic generated by the OSLF functor over a population of autonomous units factors as a product of per-unit logics; over a shared substrate it does not. The error-theoretic crossover at $r^{\ast}$ and the observability factoring under OSLF are, we conjecture, two views of the same threshold.

A further substrate the framework should illuminate is the social. Institutions, markets, languages, currencies, and standards are shared substrates with their own $\varepsilon$, $\rho$, $v$. The recurring question of when to federate versus when to centralise is, on this view, the same question as where to draw a chiplet boundary --- and should admit the same answer.

\begin{thebibliography}{9}
\bibitem{spj-triemaps} S.~Peyton~Jones, S.~Marlow, \emph{Triemaps That Match}, 2014.
\bibitem{mcbride-derivative} C.~McBride, \emph{The Derivative of a Regular Type is its Type of One-Hole Contexts}, 2001.
\bibitem{vandervorst-pathmap} A.~Vandervorst, \emph{PathMap: zipper-composable triemap implementation}.
\bibitem{flp} M.~J.~Fischer, N.~A.~Lynch, M.~S.~Paterson, \emph{Impossibility of Distributed Consensus with One Faulty Process}, JACM 1985.
\bibitem{google-sdc} P.~H.~Hochschild et al., \emph{Cores that don't count}, HotOS 2021.
\bibitem{meta-sdc} H.~D.~Dixit et al., \emph{Silent Data Corruptions at Scale}, 2021.
\bibitem{f1r3fly-oslf} F1R3FLY.io, \emph{Observer-generated logics via OSLF functors over GSLTs}, internal research note.
\bibitem{west-1997} G.~B.~West, J.~H.~Brown, B.~J.~Enquist, \emph{A General Model for the Origin of Allometric Scaling Laws in Biology}, Science 276:122--126, 1997.
\bibitem{west-book} G.~B.~West, \emph{Scale: The Universal Laws of Growth, Innovation, Sustainability, and the Pace of Life in Organisms, Cities, Economies, and Companies}, Penguin Press, 2017.
\bibitem{west-cities} L.~M.~A.~Bettencourt, J.~Lobo, D.~Helbing, C.~K\"uhnert, G.~B.~West, \emph{Growth, innovation, scaling, and the pace of life in cities}, PNAS 104:7301--7306, 2007.
\end{thebibliography}

\end{document}
